\section{\texorpdfstring{\code{driver::can::Frame}}{driver::can::Frame}}
\label{sec:framecls}

Registerkartan i \kapref{ch:can} säger exakt vad hårdvaran kan skicka och ta emot, och det är inte
mycket: en identifierare, en längd och upp till åtta databytes. Klassen speglar det:

\begin{cppcode}
namespace driver::can
{
/**
 * @brief A single CAN frame: identifier, data length code and payload.
 */
struct Frame
{
    /** Maximum number of payload bytes in a CAN frame. */
    static constexpr std::uint8_t MaxDataLen{8U};

    std::uint16_t id{};               /**< CAN identifier, 11 bits (0x000 - 0x7FF). */
    std::uint8_t  dlc{};              /**< Data length code, 0 - 8. */
    std::uint8_t  data[MaxDataLen]{}; /**< Payload. */
};

/**
 * @brief Check whether a frame is valid, i.e. whether the hardware can send it.
 */
[[nodiscard]] constexpr bool isValid(const Frame& frame) noexcept
{
    return (frame.id <= 0x7FFU) && (frame.dlc <= Frame::MaxDataLen);
}

} // namespace driver::can
\end{cppcode}

Det är hela filen. Ingen \code{.cpp}.

\subsection{Varför en aggregate}

\code{Frame} har medvetet ingen konstruktor, destruktor eller medlemsfunktion. Skälen, i tur och
ordning:

\textbf{Den är data, inte beteende.} En frame är något som flyttas: från applikationen till
drivrutinen, från drivrutinen till fyra register, från registren ut på bussen.

\textbf{Aggregate-initiering är läsbar.} Med aggregaten fungerar
\code{Frame frame{0x123U, 2U, {0xAAU, 0xBBU}}}. Lägg till en enda användardeklarerad konstruktor
och den raden slutar kompilera. Testerna du möter från och med nästa kapitel är fulla av sådana
rader.

\textbf{Default member initialization gör \code{Frame f{}} säker.} \code{{}} efter varje medlem
betyder att en default-konstruerad frame har nollställda fält. Utan dem hade \code{receive()}
kunnat lämna ifrån sig skräp i de bytes den inte fyllde.

\subsection{\texorpdfstring{\code{static constexpr}}{static constexpr}, inte \texorpdfstring{\code{\#define}}{define}}

\code{MaxDataLen} har en \textbf{typ}, en \textbf{räckvidd} och finns i \textbf{debuggern}. En
\code{#define} har inget av det, och gäller dessutom från raden den står på till filens slut, tvärs
igenom varje header som inkluderas efteråt.

Den används dessutom som arraystorlek, vilket gör att arrayens storlek och gränskontrollen i
\code{isValid()} garanterat är samma tal.

\subsection{Varför \texorpdfstring{\code{isValid()}}{isValid()} är en fri funktion}

Två skäl. Den håller aggregaten intakt --- en medlemsfunktion är första steget mot en klass med
invarianter, och nästa steg är en konstruktor. Och den hör till \emph{drivrutinen}, inte till
datan: regeln kommer från den här hårdvarans registerbredder, inte från CAN som begrepp.

\begin{aside}
\note[Regel] Hårdvaran validerar ingenting. Drivrutinen \textbf{måste} därför validera --- och
\textbf{bara} drivrutinen. Frestelsen att lägga en kontroll till i applikationen "för säkerhets
skull" ska motstås: två kopior av samma regel är två saker som ska hållas i synk, och den dag
\code{MaxDataLen} ändras kommer bara den ena att ändras med.
\end{aside}
